\chapter{Introduction}
\label{ch:Introduction}

The global energy transition is driving society toward cleaner and more sustainable ways of producing electricity. This transformation has accelerated the deployment of renewable energy sources, most of which are connected to the grid through power electronic converters, commonly referred to as Inverter Based Resources (IBRs). 
Unlike conventional synchronous machines, IBRs lack the natural inertia and damping that arise from the big rotating masses of generators. As synchronous units are progressively being decommissioned, the power system experiences reduced inertia and weaker frequency stability. This shift towards power systems with high penetration of IBRs creates new challenges in maintaining dynamic stability, especially under large disturbances.

IBRs can be broadly categorized into Grid-Following (GFL) and Grid-Forming (GFM) inverters.
Grid-Following inverters operate as controlled current sources that inject current based on an externally measured grid voltage. In contrast, Grid-Forming inverters behave as controlled voltage sources, capable of establishing their own voltage and frequency references.
This fundamental difference enables GFM inverters to provide advanced functionalities such as islanded operation, black-start capability, and voltage and frequency support during disturbances. As a result, they are seen as a promising solution to enhance the stability of low-inertia systems.
Among the technologies capable of providing GFM capability, Battery Energy Storage Systems (BESS) are particularly attractive. Their fast response and bidirectional power flow make them suitable for supporting the grid during transients while participating in energy arbitrage and ancillary service markets.
As energy markets evolve to value flexibility and stability, understanding how BESS equipped with GFM control can contribute technically and economically becomes crucial.
For RWE, many assets are connected at distribution system level, where the interaction between multiple IBRs and the grid can be complex. However, manufacturers often protect the internal control algorithms of their inverters, limiting transparency and making it challenging for developers to predict system behavior.

Therefore, it is important to understand how GFM parameters—such as virtual inertia, damping, and droop coefficients—affect system performance, and how these can be interpreted from available manufacturer data sheets to ensure appropriate equipment selection and sizing.
This thesis aims to analyze the impact of Grid-Forming parameters on system stability under large disturbances using PowerFactory simulations at the distribution network level. Different park configurations—including standalone BESS and hybrid renewable plants—will be evaluated to determine the minimum share of GFM capacity required to ensure stable operation.
The study focuses on system-level behavior and parameter sensitivity, not on designing or modifying inverter control loops. By doing so, it supports RWE in developing internal capability to assess GFM behavior and translate manufacturer information into actionable technical requirements.
Managing the stability of future power systems requires addressing the growing share of generation from nontraditional sources, including renewable technologies such as solar or wind and energy storage systems like batteries. These resources not only exhibit variability due to their weather-dependent energy inputs but also span a wide range of capacities—from small residential rooftop installations of a few kilowatts to large utility-scale plants reaching hundreds or even thousands of megawatts. Moreover, they are integrated across multiple grid levels, connecting both within distribution networks and directly to high-voltage transmission systems.

The operation of future power networks must integrate both the inherent dynamics and control responses of traditional large-scale synchronous machines and those of emerging inverter-based generation technologies.

The thesis provides a quantitative estimate of the minimum GFM share required for acceptable dynamic performance across different grid strengths, it also provide  

% ...



%Am Ende des Einführungskapitels ein kurzer Überblick über die kommenden Kapitel in einem Absatz (Beschreibung der Struktur der Thesis):

%...